\documentclass[11pt,a4paper]{article}
\usepackage[a4paper,margin=2.5cm]{geometry}
\usepackage{amsmath,amssymb,amsthm}
\usepackage{booktabs}
\usepackage[colorlinks=true,linkcolor=blue!60!black]{hyperref}

\newtheorem{prop}{Proposition}
\newtheorem{cor}{Corollary}
\theoremstyle{definition}
\newtheorem{meas}{Measured fact}
\theoremstyle{remark}
\newtheorem{rem}{Remark}

\title{\textbf{Why the Hypotheses Failed:\\ The Saturation Obstruction in the Softmax--Fisher Geometry}\\[1ex]
\large Post-mortem of the experimental cycle on selective SSMs}
\author{InfoRank working notes}
\date{\today}

\begin{document}
\maketitle

\section{Setup and standing notation}

We use the framework of the proposal. A selective SSM produces pre-head states
$h_t \in \mathbb{R}^d$; the readout is
\begin{equation}
p(h) \;=\; \mathrm{softmax}(Mh+b), \qquad M \in \mathbb{R}^{V\times d},
\label{eq:readout}
\end{equation}
inducing the regularized Fisher metric on $\mathbb{R}^d$,
\begin{equation}
g_\varepsilon(h) \;=\; M^\top \Sigma\!\left(p(h)\right) M + \varepsilon I,
\qquad
\Sigma(p) \;=\; \operatorname{diag}(p) - pp^\top .
\label{eq:metric}
\end{equation}
$\mathcal{R}_\varepsilon(h)$ denotes the scalar curvature of $(\mathbb{R}^d, g_\varepsilon)$,
$\mathrm{erank}(h) = \exp H\!\big(\lambda/\operatorname{tr}\big)$ the effective rank of the spectrum,
and $\Delta_t \in \mathbb{R}^d_+$ the input-dependent discretization of the SSM block. Crucially,
$\Delta_t$ enters only the \emph{dynamics} $h_{t+1} = \Phi(h_t; x_{t+1}, \Delta_{t+1}, \theta)$ ---
never the readout \eqref{eq:readout}.

All empirical claims cited below are from the reported cycle (copy task, $d = 64$, $V = 64$, three to
five seeds; $\varepsilon = 10^{-3}\,\operatorname{tr}(g)/d$ unless stated otherwise).

The three failures share one root cause:

\begin{quote}
\emph{The Fisher geometry \eqref{eq:metric} is a functional of the output distribution $p(h)$
alone, and it degenerates --- becomes flat and isotropic --- exactly where the readout saturates
($p \to \delta$, i.e.\ $H(p) \to 0$), because $\Sigma(p) \to 0$ there. Capacity loss in the SSM
turned out to be a migration of the trajectory into precisely this saturated regime
(``confident garbage''). Every hypothesis in the proposal attached its observable to the wrong end
of the concentration scale.}
\end{quote}

\section{RQ1: curvature does not see the selection parameter}

\subsection{Obstruction 1: the factorization --- there is no channel from $\Delta$ to $\mathcal{R}$}

\begin{prop}[State-function factorization]
$\mathcal{R}_\varepsilon(h) = \mathcal{F}\big(\mathrm{softmax}(Mh+b)\big)$ for a functional
$\mathcal{F}$ of the distribution alone. In particular $\mathcal{R}_\varepsilon(h_t)$ is a function
of $h_t$ only, and
\[
\rho\big(\mathcal{R}_\varepsilon(h_t),\, \Delta_t\big)
\;=\;
\operatorname{Corr}\Big(\mathcal{F}\!\big(p(h_t)\big),\, \Delta_t\Big),
\]
so any correlation is mediated \emph{entirely} by the joint law of the trajectory $(h_t, \Delta_t)$
generated by the data distribution --- never by the geometry itself.
\end{prop}

\begin{proof}
Immediate from \eqref{eq:metric}: neither $M$, $b$ nor $h_t$ involve $\Delta_t$; the identity holds
by the law of total correlation.
\end{proof}

RQ1 therefore requires a \emph{statistical} coupling: reset positions (large $\Delta_t$) must push
the state into high-curvature regions. No mechanism enforces this. In the copy task the state
distribution at almost all positions is the stationary response to noise tokens; memory is carried
by a low-dimensional slow subspace whose decay rate $e^{-\Delta_t A}$ --- a dynamical gain ---
leaves no invariant trace in the readout geometry except through where the trajectory happens to
land.

\begin{meas}
$\rho(\mathcal{R}_\varepsilon(h_t), \Delta_t) \approx 0$ in $4/6$ models ($|\rho| < 0.05$, n.s.),
$-0.21$ in one, $+0.04$ in another; the lag profile is flat ($|\rho(\mathrm{lag})| \le 0.21$ for
$\mathrm{lag} \in [-3,3]$). Under an $\varepsilon$-ablation over two decades, $\rho$ changes sign
($-0.012 / -0.115 / +0.021$ and $+0.044 / +0.105 / -0.196$): whatever coupling exists is
$\varepsilon$-conditional, hence not a property of the geometry.
\end{meas}

\subsection{Obstruction 2: the flat limit is saturation}

The intuition behind RQ1 --- ``memory resets are violent events, so the manifold bends sharply
there'' --- conflates two different objects: the magnitude of the state update (a tangent vector)
with the bending of the readout map (a second-derivative invariant of $p$'s shape). The geometry
has the opposite behavior at the extremes of confidence:

\begin{prop}[Exponential flatness under saturation]
Write $p = (1-\delta)\,e_v + \delta\,q$ with $q$ on the remaining $V-1$ tokens, $\delta \to 0$.
Then for fixed $\varepsilon > 0$ and fixed $M$:
\[
\Sigma(p) = O(\delta), \qquad
\partial_h \Sigma = O(\delta), \qquad
\partial_h^2 \Sigma = O(\delta),
\]
hence
$g_\varepsilon = \varepsilon I + O(\delta)$,
$\Gamma = O(\delta/\varepsilon)$, and
$\mathcal{R}_\varepsilon = O(\delta/\varepsilon) \to 0$.
\end{prop}

\begin{proof}
Each softmax derivative carries a factor $p$: $\nabla_h p = J_u M$ with
$J_u = \operatorname{diag}(p) - pp^\top = O(\delta)$, and likewise $\nabla_h^2 p = O(\delta)$
(all entries contain products of two softmax factors). This gives the three displayed bounds via
\eqref{eq:metric}. The curvature formula is polynomial in $g_\varepsilon^{-1}\partial g_\varepsilon$
and $g_\varepsilon^{-1}\partial^2 g_\varepsilon$; with
$g_\varepsilon^{-1} = \tfrac1\varepsilon(I + O(\delta/\varepsilon))$ the claim follows.
\end{proof}

\begin{rem}
Two consequences worth spelling out.
\emph{(a) Saturated regions are geometrically invisible:} a one-hot readout has zero Fisher
covariance ($\Sigma(\delta_v) = 0$ identically), so ``confident garbage'' is flat, not curved.
Any expectation that resets or capacity breakdown produce curvature spikes presumes that
high-confidence or high-magnitude events bend the manifold; they do not --- they \emph{switch it
off}.
\emph{(b) The residual curvature of nearly-saturated regions is $\varepsilon$-amplified:}
$\mathcal{R}_\varepsilon = O(\delta/\varepsilon)$ grows as the floor $\varepsilon$ is lowered.
This matches the measured scaling of $|\mathcal{R}|$ (mean $-197 \to -7300$ for a $10\times$
smaller $\varepsilon$) and is why only scale-free conclusions (drops, sign-stability across
$\varepsilon$, correlations at fixed $\varepsilon$) are admissible.
\end{rem}

\begin{meas}
$\rho(\mathcal{R}_\varepsilon(h_t), \|h_t - h_{t-1}\|) < 0$ in all base models ($-0.28$ to
$-0.66$): curvature peaks where the trajectory moves \emph{slowly} --- the state parks in stiff
regions of the partial-concentration band --- and is uninformative in the saturated and relaxed
regimes that dominate reset positions.
\end{meas}

\section{RQ2: scale--shape conflation --- why $\mathrm{erank}$ moved the wrong way}

\subsection{Obstruction 3: $\mathrm{erank}$ is scale-invariant, information is not}

\begin{prop}[Scale invariance]
For $c > 0$: $\mathrm{erank}(c\,g) = \mathrm{erank}(g)$, with $\mathrm{erank} = d$ if and only if
the normalized spectrum is uniform. Hence $\mathrm{erank}$ measures only the \emph{shape} of the
spectrum; all information in the \emph{scale} ($\operatorname{tr} g$, $\det g$) is invisible to it.
\end{prop}

\subsection{What actually happens beyond capacity}

\begin{meas}
For models trained at $L = 257$ and evaluated OOD at $L \in \{513, 1025\}$: accuracy falls to
chance; the per-position output entropy collapses to $H(p) < 0.001$ nats (numerically one-hot
outputs --- \emph{confident garbage}, not uncertainty); the cross-sequence state dispersion falls
by ${\sim}4\times$ ($1.33 \to 0.33$): the state collapses onto a low-variance attractor; and the
metric degenerates: $\mathrm{cond}(g_\varepsilon) \to 1.0000$ per point,
$\overline{\log\det} \approx -573$ (vs.\ $-419$ in capacity), $\mathrm{erank} \to d = 64$.
\end{meas}

This is Proposition 2 enacted by the training dynamics, on two levels:
\begin{enumerate}
\item \emph{Dynamical level:} beyond capacity the retained signal decays below the noise floor;
the loss-optimal behavior is a fixed attractor response, so the trajectory loses both memory
content and input-driven fluctuation (dispersion ${\sim}4\times$ down).
\item \emph{Geometric level:} the readout saturates ($H(p) \to 0$). By Proposition 2,
$g_\varepsilon \to \varepsilon I + O(\delta)$: the metric collapses onto the isotropic
$\varepsilon$-floor. Points with $O(\delta) \lesssim \varepsilon$ sit exactly on the floor
(eigenvalue $\varepsilon$, condition number $1$); the ensemble mean $\bar\lambda \approx
1.3\times10^{-4} = 9\,\varepsilon$ reflects the admixture of not-fully-saturated positions.
\end{enumerate}

\begin{cor}[The RQ2 inversion]
Along this path, the in-band anisotropy ($\mathrm{cond} \sim 10^{3.5\text{--}3.9}$,
$\mathrm{erank} \approx 15$) gives way to a flat spectrum ($\mathrm{cond} \to 1$,
$\mathrm{erank} \to d$). Because $\mathrm{erank}$ is scale-invariant (Proposition 3), it records
\emph{only} the flattening: $\mathrm{erank}$ \emph{increases} with $L$ (fitted slope
$+0.002\ldots+0.055$, $R^2 = 0.51$--$0.91$) and correlates \emph{negatively} with accuracy
($r$ up to $-0.98$). The hypothesized law $\mathrm{erank}(h_L) = \alpha - \beta L$ with
$\beta > 0$ cannot hold in a saturating regime. The correct signature is
\[
\text{capacity loss}
\;\Longleftrightarrow\;
\operatorname{tr}(g_\varepsilon)\!\downarrow,\;\;
\det(g_\varepsilon)\!\downarrow,\;\;
\mathrm{cond}(g_\varepsilon) \to 1,
\qquad
\mathrm{erank} \uparrow d \;\;\text{(anti-indicator)}.
\]
\end{cor}

\begin{rem}
The saturation mechanism also resolves an apparent puzzle: the output becomes maximally
\emph{certain} ($H(p) \to 0$) while the latent metric becomes maximally \emph{flat}
($\mathrm{erank} \to d$). Fisher information about $h$ vanishes not because the model is unsure,
but because it is sure --- $\Sigma(p) \to 0$ kills the metric regardless of whether the confidence
is warranted. Scale statistics can tell warranted from unwarranted confidence (they fall only in
the unwarranted case, relative to the in-capacity baseline); $\mathrm{erank}$ cannot.
\end{rem}

\section{RQ3: the confidence--isotropy tension}

The penalty is
$\mathcal{L}_{\mathrm{reg}} = \gamma\,\mathbb{E}_t\big\|\log g_\varepsilon(h_t) - \log(\tau
I)\big\|_F^2 = \gamma\,\mathbb{E}_t\sum_i (\log \lambda_i - \log \tau)^2$.

\begin{prop}[What the penalty optimizes]
\mbox{}\phantom{}
\begin{enumerate}
\item[(i)] The per-eigenvalue gradient $2(\log \lambda_i - \log\tau)/\lambda_i$ drives every
$\lambda_i \to \tau$: the minimizer over metrics is $g = \tau I$, the flat Euclidean metric.
\item[(ii)] Since $\Sigma(p) \preceq I$ (as $a^\top \Sigma a = \operatorname{Var}_p(a_v) \le
\|a\|^2$), we have $\lambda_i(g_\varepsilon) \le \lambda_i(M^\top M) + \varepsilon$: the penalty
can raise the $\lambda_i$ only by raising the spread of $p$ away from one-hot. It is a \emph{reward
for output uncertainty}.
\item[(iii)] The task loss acts oppositely: on learnable positions CE drives $p$ toward one-hot,
whence $\Sigma(p) \to 0$ and $\lambda_i \to \varepsilon$.
\end{enumerate}
\end{prop}

Three consequences, all observed:
\begin{enumerate}
\item \emph{Adversarial pressure on CE.} The objectives fight on every confident prediction. At
$\gamma = 0.1$ this appeared as training instability (one of four regularized seeds collapsed to
$4.5\%$ accuracy); the surviving seeds match baseline.
\item \emph{No benefit.} Final accuracy $18.3\%$ (reg, $n{=}4$) vs.\ $22.1\%$ (base, $n{=}3$),
Welch $p = 0.52$. The hypothesized mechanism --- ``the barrier prevents spectral collapse, hence
preserves recall'' --- misidentifies the failure mode: by Section 3, long-context breakdown is
\emph{saturation} ($\lambda_i \to \varepsilon$ from the confident end, with $H(p) \to 0$), not
uncertainty-driven collapse; and the in-band anisotropy that the penalty suppresses is exactly
what carries task information (Corollary 1).
\item \emph{Self-undermining diagnostics.} At the penalty's minimizer $g = \tau I$ the manifold is
flat: no curvature, no phase structure, nothing left to measure. The regularizer eliminates the
very signal this research program wants to track. The observed ``metric stabilization'' (erank
variance across a $16\times$ length range: $\pm 0.8$ regularized vs.\ $\pm 15$ baseline) is bought
by destroying the metric's information content --- metric stability is not capability.
\end{enumerate}

\section{Why the one robust signal survives}

The single finding that replicated with decisive statistics --- the erank drop $44 \to 16$ at the
write/read boundary (bootstrap CIs excluding zero by orders of magnitude; $\varepsilon$-invariant
over two decades: $27.65/27.70/27.63$ and $19.20/19.24/19.18$) --- survives because it lives in the
partial-concentration band and is \emph{phase-locked to a task event}: during writing, predictions
across positions cover the full token inventory with mixed confidence ($\Sigma$ of high effective
dimension); during readout they concentrate ($\Sigma$ low-rank). It is a systematic transition of
the \emph{shape} of $p$ within the band --- not a correlation with a dynamical gain (RQ1's error),
not a scale confound (RQ2's error), and not an optimization target (RQ3's error).

\section{Revised observables}

\begin{center}
\begin{tabular}{@{}llp{6.9cm}@{}}
\toprule
\textbf{Replace} & \textbf{with} & \textbf{reason} \\
\midrule
$\rho(\mathcal{R}_\varepsilon(h_t), \Delta_t)$ &
event-aligned $\mathrm{erank}(h_t)$ &
$\mathcal{R}$ has no channel from the dynamics (Prop.\ 1) and vanishes exponentially under
saturation (Prop.\ 2); regime switches are visible in the spectrum's shape, phase-locked to task
events. \\
$\mathrm{erank}(h_L) = \alpha - \beta L$ &
$\log\det g_\varepsilon$, $\mathrm{cond}\,g_\varepsilon$ vs.\ $L$ &
$\mathrm{erank}$ is scale-invariant (Prop.\ 3) and \emph{rises} under the saturation collapse;
scale statistics fall monotonically and track accuracy ($r$ up to $0.95$). \\
$\mathcal{L}_{\mathrm{reg}}$ on $\log \lambda_i$ &
penalize $\mathrm{cond}$, or drop it &
the log-eigenvalue penalty rewards uncertainty, opposes CE (Prop.\ 4), and its minimizer destroys
the diagnostic geometry. \\
\bottomrule
\end{tabular}
\end{center}

\end{document}
